% !TEX root = ../main.tex
\corrigesChapitre{12}{Description d'un fluide au repos}
\corrigesSouples

\begin{blocLivre}{cor-12-1}
\begin{corrige}{1}{Du microscopique au macroscopique}
\demarche{La température traduit l'agitation des molécules, la pression leurs chocs sur les parois ; la masse vaut $m=\rho\,V$.}
\begin{questions}
\item Plus la température est élevée, plus les molécules sont \res{\text{rapides}} (schéma ci-dessous).
\item Elles frappent les parois \res{\text{plus souvent et plus fort}} : la pression \res{\text{augmente}}.
\item $V=5\times4\times2{,}5=\qty{50}{m^3}$, d'où $m=\rho\,V=1{,}2\times50=\res{\qty{60}{kg}}$.
\item Molécules très \res{\text{éloignées}} dans un gaz, au \res{\text{contact}} dans un liquide : bien plus de matière par litre.
\end{questions}
\schema[corrige]{12-ex01-cor}
\retiens{À volume constant, chauffer un gaz augmente sa pression : on contrôle les pneus « à froid ».}
\end{corrige}
\end{blocLivre}
\begin{blocLivre}{cor-12-2}
\begin{corrige}{2}{La force de l'air sur une fenêtre}
\demarche{$F=P\times S$ avec $P$ en Pa et $S$ en \unit{m^2} ; la vitre subit deux forces pressantes opposées, l'une de chaque côté.}
\begin{questions}
\item $S=1{,}20\times0{,}80=\qty{0,96}{m^2}$ et $F=\num{1,013e5}\times0{,}96=\res{\qty{9,72e4}{N}}$.
\item $m=\dfrac{F}{g}=\dfrac{\num{9,72e4}}{9{,}81}\approx\res{\qty{9,9e3}{kg}}$, soit près de \res{\text{10 tonnes}}.
\item L'air intérieur exerce une force \res{\text{opposée}}, de même valeur : les deux forces se compensent.
\item $\Delta P=\qty{30}{hPa}=\qty{3000}{Pa}$ : $F=3000\times0{,}96=\res{\qty{2,88e3}{N}}$, vers l'\res{\text{extérieur}}.
\end{questions}
\schema[corrige]{12-ex02-cor}
\piege{Oublier de convertir : $\qty{1}{hPa}=\qty{100}{Pa}$. Ici, \qty{30}{hPa} suffisent à pousser la vitre avec près de \qty{3000}{N}.}
\end{corrige}
\end{blocLivre}
\begin{blocLivre}{cor-12-3}
\begin{corrige}{3}{La seringue bouchée}
\demarche{Loi de Mariotte $P_1\,V_1=P_2\,V_2$ pour l'air enfermé ; le piston immobile est soumis à des forces qui se compensent.}
\begin{questions}
\item La \res{\text{loi de Mariotte}} : la quantité d'air est fixée (seringue bouchée) et la température reste constante (compression lente).
\item $P_2=\dfrac{P_1\,V_1}{V_2}=\dfrac{1013\times60}{20}=\res{\qty{3039}{hPa}}=\qty{3,04e5}{Pa}$.
\item Piston en équilibre : $F_\mathrm{op}+P_\mathrm{atm}\,S=P_2\,S$, d'où $F_\mathrm{op}=(\num{3,039e5}-\num{1,013e5})\times\num{6,0e-4}\approx\res{\qty{122}{N}}$.
\item $V=\dfrac{P_1\,V_1}{P}=\dfrac{\num{1,013e5}\times60}{\num{2,0e5}}\approx\res{\qty{30,4}{mL}}$.
\end{questions}
\schema[corrige]{12-ex03}
\retiens{Dans $P_1\,V_1=P_2\,V_2$, on peut garder n'importe quelles unités (hPa, mL…), pourvu que ce soient les mêmes des deux côtés. Dans
$F=P\times S$, en revanche, le Pa et le \unit{m^2} sont obligatoires : $\qty{6,0}{cm^2}=\qty{6,0e-4}{m^2}$.}
\end{corrige}
\end{blocLivre}
\begin{blocLivre}{cor-12-4}
\begin{corrige}{4}{La pression sous l'eau}
\demarche{À la profondeur $h$ : $P=P_\mathrm{atm}+\rho\,g\,h$ (loi de la statique entre la surface et le point).}
\begin{questions}
\item $P=\num{1,013e5}+\num{1,00e3}\times9{,}81\times2{,}0=\num{1,013e5}+\num{1,96e4}\approx\res{\qty{1,21e5}{Pa}}$.
\item $P=\num{1,013e5}+\num{1,03e3}\times9{,}81\times10\approx\res{\qty{2,02e5}{Pa}}$, soit environ $2\,P_\mathrm{atm}$.
\item $h=\dfrac{P-P_\mathrm{atm}}{\rho\,g}=\dfrac{\num{3,0e5}-\num{1,013e5}}{\num{1,03e3}\times9{,}81}\approx\res{\qty{19,7}{m}}$.
\item Une \res{\text{fonction affine}} : une droite d'ordonnée à l'origine $P_\mathrm{atm}$.
\end{questions}
\schema[corrige]{12-ex04-cor}
\end{corrige}
\end{blocLivre}
\begin{blocLivre}{cor-12-5}
\begin{corrige}{5}{Des vases communicants}
\demarche{Dans un liquide au repos, la pression augmente avec la profondeur et ne dépend que de l'altitude du point.}
\begin{questions}
\item Du moins profond au plus profond : $\res{P_\mathrm{D}<P_\mathrm{A}<P_\mathrm{E}<P_\mathrm{C}}$.
\item A et B sont au même niveau dans la même eau au repos : $\res{P_\mathrm{A}=P_\mathrm{B}}$.
\item $P_\mathrm{A}=P_\mathrm{atm}+\rho\,g\,(z_\mathrm{D}-z_\mathrm{A})=\num{1,013e5}+\num{1,00e3}\times9{,}81\times1{,}0\approx\res{\qty{1,11e5}{Pa}}$ ;
$P_\mathrm{C}=\num{1,013e5}+\num{1,00e3}\times9{,}81\times2{,}7\approx\res{\qty{1,28e5}{Pa}}$.
\item Les deux surfaces libres sont à la pression $P_\mathrm{atm}$ : deux points de même pression d'un même liquide au repos sont
au \res{\text{même niveau}}.
\end{questions}
\schema[corrige]{12-ex05}
\retiens{La pression ne dépend ni de la forme du récipient ni de la quantité de liquide au-dessus : seulement de la profondeur.}
\end{corrige}
\end{blocLivre}
\begin{blocLivre}{cor-12-6}
\begin{corrige}{6}{La pression en plongée, en Python}
\demarche{Le programme calcule $P=P_0+\rho\,g\,h$ tous les \qty{10}{m} et l'affiche en bar ($\qty{1}{bar}=\qty{1e5}{Pa}$).}
\begin{questions}
\item La loi de la statique $P=P_0+\rho\,g\,h$ ; \code{P0} est la pression \res{\text{atmosphérique}}, à la surface.
\item \code{0 m : 1.01 bar}, \code{10 m : 2.02 bar}, \code{20 m : 3.03 bar}, \code{30 m : 4.04 bar}, \code{40 m : 5.05 bar}.
\item $\rho\,g\times10=\qty{1,01e5}{Pa}$ : \res{\qty{1,01}{bar}} tous les \qty{10}{m}. Triple : $\rho\,g\,h=2\,P_0$, d'où $h=\dfrac{2\times\num{1,013e5}}{1030\times9{,}81}\approx\res{\qty{20,1}{m}}$.
\item Première ligne : \code{rho, g, P0 = 1000, 9.81, 0.80e5}. À \qty{20}{m} : $\num{0,80e5}+1000\times9{,}81\times20\approx\res{\qty{2,76}{bar}}$.
\end{questions}
\end{corrige}
\end{blocLivre}
\begin{blocLivre}{cor-12-7}
\begin{corrige}{7}{Les poumons de l'apnéiste}
\demarche{Pression par la loi de la statique, puis loi de Mariotte pour l'air des poumons, à température constante.}
\begin{questions}
\item $P=\num{1,013e5}+\num{1,03e3}\times9{,}81\times20\approx\res{\qty{3,03e5}{Pa}}$.
\item $V=\dfrac{P_0\,V_0}{P}=\dfrac{\num{1,013e5}\times6{,}0}{\num{3,03e5}}\approx\res{\qty{2,0}{L}}$ : le volume est divisé par trois.
\item Le volume décroît de moins en moins vite (courbe ci-dessous) : \qty{3,0}{L} à \qty{10}{m}, \qty{1,5}{L} à \qty{30}{m}.
\item \qty{6,0}{L} d'air respiré à \qty{3,03e5}{Pa} occuperaient en surface $\dfrac{6{,}0\times\num{3,03e5}}{\num{1,013e5}}\approx\res{\qty{18}{L}}$ : les poumons se
\res{\text{déchireraient}}. Le plongeur doit expirer en remontant.
\end{questions}
\schema[corrige]{12-ex07-cor}
\piege{Prendre $\rho\,g\,h$ pour la pression : Mariotte utilise la pression totale $P_\mathrm{atm}+\rho\,g\,h$.}
\end{corrige}
\end{blocLivre}
\begin{blocLivre}{cor-12-8}
\begin{corrige}{8}{Le baromètre de Torricelli}
\demarche{A est au niveau de la surface libre de la cuve ; au-dessus de la colonne règne le vide.}
\begin{questions}
\item $\res{P_\mathrm{C}=0}$ (vide) ; A est au même niveau que le point B de la surface de la cuve : $\res{P_\mathrm{A}=P_\mathrm{atm}}$.
\item $P_\mathrm{A}-P_\mathrm{C}=\rho\,g\,(z_\mathrm{C}-z_\mathrm{A})=\rho\,g\,h$, d'où $P_\mathrm{atm}=\num{13,6e3}\times9{,}81\times0{,}760\approx\res{\qty{1,01e5}{Pa}}$.
\item La hauteur est \res{\text{proportionnelle}} à $P_\mathrm{atm}$ : \qty{760}{mm} de mercure correspondent à \qty{1013}{hPa} (tension artérielle).
\item $h=P_\mathrm{atm}/(\rho_\mathrm{eau}\,g)\approx\res{\qty{10,3}{m}}$ : impraticable ; le mercure, 13,6 fois plus dense, donne une colonne 13,6 fois plus courte.
\end{questions}
\schema[corrige]{12-ex08}
\end{corrige}
\end{blocLivre}
\begin{blocLivre}{cor-12-9}
\begin{corrige}{9}{La pression et l'altitude}
\demarche{Avec $\rho$ constante, la pression diminue de $\rho\,g\,\Delta z$ quand on s'élève de $\Delta z$.}
\begin{questions}
\item $1{,}2\times9{,}81\times100\approx\res{\qty{1,18e3}{Pa}}$, soit environ \qty{12}{hPa} ; pour \qty{300}{m} : environ \res{\qty{35}{hPa}}.
\item Pour un étage : $1{,}2\times9{,}81\times3\approx\qty{35}{Pa}=\qty{0,35}{hPa}$, plus que \qty{0,1}{hPa} : le téléphone repère \res{\text{chaque étage}}.
\item $1013-\dfrac{1{,}2\times9{,}81\times4807}{100}\approx1013-566=\res{\qty{447}{hPa}}$.
\item L'air est \res{\text{compressible}} : sa masse volumique diminue avec l'altitude, et la pression décroît de moins en moins vite.
L'hypothèse « $\rho$ constante » n'est valable que sur quelques centaines de mètres.
\end{questions}
\retiens{La loi $P_2-P_1=\rho\,g\,(z_1-z_2)$ suppose un fluide incompressible, de masse volumique constante : parfaite pour les liquides,
acceptable pour l'air seulement sur de faibles dénivelés.}
\end{corrige}
\end{blocLivre}
\begin{blocLivre}{cor-12-10}
\begin{corrige}{10}{Mesurer une masse volumique avec un tube en U}
\demarche{On égale les pressions en deux points de même niveau d'un même liquide continu : ici l'eau, en A et B.}
\begin{questions}
\item A et B sont dans la \res{\text{même eau}}, continue et au repos, au même niveau : ils ont la même pression.
\item $P_\mathrm{A}=P_\mathrm{atm}+\rho_\mathrm{h}\,g\,h_\mathrm{h}$ et $P_\mathrm{B}=P_\mathrm{atm}+\rho_\mathrm{eau}\,g\,h_\mathrm{e}$.
\item $\rho_\mathrm{h}\,h_\mathrm{h}=\rho_\mathrm{eau}\,h_\mathrm{e}$, d'où $\rho_\mathrm{h}=\dfrac{\num{1,00e3}\times10{,}8}{12{,}0}=\res{\qty{900}{kg.m^{-3}}}$.
\item $P_\mathrm{A}=\num{1,013e5}+900\times9{,}81\times0{,}120\approx\res{\qty{1,024e5}{Pa}}$.
\end{questions}
\schema[corrige]{12-ex10}
\piege{Égaler les pressions en deux points situés dans des liquides différents : au niveau de la surface de l'eau, par exemple, le
point de la branche droite est dans l'huile, et les pressions ne sont pas égales.}
\end{corrige}
\end{blocLivre}
\begin{blocLivre}{cor-12-11}
\begin{corrige}{11}{Tester la loi de Mariotte}
\demarche{Si $P\times V$ est constant, $P=\dfrac{k}{V}$ : le graphe de $P$ en fonction de $\dfrac{1}{V}$ est une droite passant par l'origine.}
\begin{questions}
\item Produits en \unit{hPa.mL} : \num{60720} ; \num{60900} ; \num{60600} ; \num{60930} ; \num{60624} ; \num{60920}. Ils valent tous \res{\num{6,08e4}} à \qty{0,3}{\percent} près : la loi
est \res{\text{vérifiée}}.
\item Une \res{\text{droite passant par l'origine}} (graphe ci-dessous).
\item $P=\dfrac{\num{6,08e4}}{15}\approx\res{\qty{4,05e3}{hPa}}$.
\item Une compression rapide \res{\text{échauffe}} l'air : la température ne serait plus constante. L'air du \res{\text{tuyau}} reliant la seringue au
capteur n'est pas compté dans $V$ : il faut un tuyau très court.
\end{questions}
\schema[corrige]{12-ex11-cor}
\end{corrige}
\end{blocLivre}
\begin{blocLivre}{cor-12-12}
\begin{corrige}{12}{La plongée en bouteille}
\demarche{Statique pour la pression à \qty{20}{m}, puis loi de Mariotte pour ramener tous les volumes d'air à \qty{1,0}{bar}.}
\begin{questions}
\item $P=\num{1,0e5}+\num{1,03e3}\times9{,}81\times20\approx\qty{3,02e5}{Pa}=\res{\qty{3,02}{bar}}$.
\item $V=\dfrac{200\times12}{1{,}0}=\res{\qty{2400}{L}}$ ; avant la réserve, $\dfrac{(200-50)\times12}{1{,}0}=\res{\qty{1800}{L}}$ sont utilisables.
\item \qty{20}{L} à \qty{3,02}{bar} occupent $\dfrac{20\times3{,}02}{1{,}0}=\res{\qty{60,4}{L}}$ à la pression atmosphérique, chaque minute.
\item $\dfrac{1800}{60{,}4}\approx\res{\qty{30}{min}}$ à \qty{20}{m}, contre $\dfrac{1800}{20}=\res{\qty{90}{min}}$ près de la surface : l'autonomie est divisée par trois.
\end{questions}
\retiens{En profondeur, chaque inspiration contient plus de molécules d'air, car l'air respiré est comprimé : la bouteille se vide plus
vite.}
\end{corrige}
\end{blocLivre}
\begin{blocLivre}{cor-12-13}
\begin{corrige}{13}{Le ballon-sonde}
\demarche{Loi de Mariotte pour l'hélium, puis $r=\left(\dfrac{3\,V}{4\,\pi}\right)^{1/3}$ pour le rayon.}
\begin{questions}
\item $V_2=\dfrac{1013\times1{,}5}{30}\approx\res{\qty{50,7}{m^3}}$ : le volume est multiplié par environ \res{34}.
\item Au sol, $r=\left(\dfrac{3\times1{,}5}{4\pi}\right)^{1/3}\approx\qty{0,71}{m}$ : $\res{d\approx\qty{1,4}{m}}$ ; à \qty{24}{km}, $r\approx\qty{2,30}{m}$ : $\res{d\approx\qty{4,6}{m}}$.
\item $V_\mathrm{max}=\dfrac{4}{3}\,\pi\times3{,}0^3\approx\qty{113}{m^3}$, d'où $P=\dfrac{1013\times1{,}5}{113}\approx\res{\qty{13,4}{hPa}}$.
\item En se refroidissant, le gaz occupe un volume \res{\text{plus petit}} que prévu : l'hypothèse « température constante » est
fausse. Le ballon éclate donc à une pression plus faible, c'est-à-dire plus haut.
\end{questions}
\schema[corrige]{12-ex13-cor}
\end{corrige}
\end{blocLivre}
\begin{blocLivre}{cor-12-14}
\begin{corrige}{14}{Le château d'eau}
\demarche{Eau au repos : surpression $\rho\,g\,h$, où $h$ est la hauteur d'eau au-dessus du point.}
\begin{questions}
\item $h=52-2{,}0=\qty{50}{m}$ : $\Delta P=\num{1,00e3}\times9{,}81\times50\approx\res{\qty{4,9e5}{Pa}}$, soit \qty{4,9}{bar}.
\item $h=52-25=\qty{27}{m}$ : $\Delta P\approx\res{\qty{2,65e5}{Pa}}$, soit \qty{2,65}{bar} $<\qty{3,0}{bar}$ : le confort \res{\text{n'est pas assuré}}.
\item $h=\dfrac{\num{3,0e5}}{\num{1,00e3}\times9{,}81}\approx\qty{30,6}{m}$ : la surface de l'eau devrait être à $25+30{,}6\approx\res{\qty{55,6}{m}}$ d'altitude.
\item Au fond de la cuve, sous \qty{8,0}{m} d'eau : $\Delta P=\qty{7,85e4}{Pa}$ et $F=\num{7,85e4}\times0{,}50\approx\res{\qty{3,92e4}{N}}$.
\end{questions}
\schema[corrige]{12-ex14}
\retiens{Seule compte la hauteur d'eau : les châteaux d'eau sont bâtis sur les points hauts.}
\end{corrige}
\end{blocLivre}
\begin{blocLivre}{cor-12-15}
\begin{corrige}{15}{Pourquoi les barrages sont-ils si épais à la base ?}
\demarche{La surpression due à l'eau vaut $\rho\,g\,h$ : elle est proportionnelle à la profondeur, donc sa moyenne sur la paroi vaut la moitié
de sa valeur au pied.}
\begin{questions}
\item À \qty{30}{m} : $\num{1,00e3}\times9{,}81\times30\approx\res{\qty{2,94e5}{Pa}}$ ; au pied, à \qty{60}{m} : $\res{\qty{5,89e5}{Pa}}$.
\item Des forces \res{\text{perpendiculaires}} à la paroi, dirigées de l'eau vers le béton, de plus en plus grandes vers le bas (schéma).
\item $F=\num{5,89e5}\times1{,}0=\res{\qty{5,89e5}{N}}$.
\item Pression moyenne \qty{2,94e5}{Pa} sur $60\times1=\qty{60}{m^2}$ : $F\approx\res{\qty{1,77e7}{N}}$ par mètre de largeur, le poids de \num{1800} tonnes. Les
efforts étant maximaux en bas, le barrage doit y être \res{\text{très épais}}.
\end{questions}
\schema[corrige]{12-ex15}
\end{corrige}
\end{blocLivre}
